\chapter{Classical Randomized Algorithms: A Guided Core}
\label{ch:stage3-classical-randomized}

This supplement gathers the classical algorithms that students should master
before moving to specialized applications.  The book already develops
randomized selection, game-tree evaluation, complexity classes, and
amplification in other chapters; those topics remain in place and are
cross-referenced here.

\section{Las Vegas and Monte Carlo algorithms}

\begin{definition}
A Las Vegas algorithm always returns a correct answer and has a random
running time.  A Monte Carlo algorithm has a bounded running time but may
return an incorrect answer with a stated probability.
\end{definition}

For a Las Vegas algorithm, the main guarantee is usually
$\mathbb{E}[T]\leq t(n)$, together with a tail bound when available.  For a
Monte Carlo algorithm, the statement must include whether the error is one-
sided or two-sided and whether the probability is over the algorithm's
randomness for every fixed input.

\begin{proposition}[Truncation]
If a Las Vegas algorithm has expected running time at most $T$ on every
input, stopping it after $T/\delta$ steps and returning ``failure'' gives a
Monte Carlo algorithm that fails with probability at most $\delta$.
\end{proposition}

\begin{proof}
By Markov's inequality, $\Pr[\text{runtime}>T/\delta]\leq\delta$.
\end{proof}

\section{Randomized Quicksort}

Randomized Quicksort chooses a uniformly random pivot, partitions around it,
and recursively sorts the two subarrays.

\begin{theorem}[Expected comparisons]
For $n$ distinct keys, randomized Quicksort performs at most
$2(n+1)H_n-4n$ expected comparisons, where
$H_n=\sum_{i=1}^{n}1/i$.  In particular, the expected cost is
$O(n\log n)$.
\end{theorem}

\begin{proof}
Let $X_{ij}$ indicate that the keys of ranks $i<j$ are compared.  They are
compared exactly when the first pivot chosen from ranks $i,\ldots,j$ is one
of the endpoints.  That event has probability $2/(j-i+1)$.  Therefore
\[
 \mathbb{E}[C_n]=\sum_{1\leq i<j\leq n}\frac{2}{j-i+1}
 =2(n+1)H_n-4n.
\]
\end{proof}

The same argument gives a clean student lesson: the subproblem sizes are
dependent, but indicator variables and linearity of expectation avoid the
need to analyze the entire recursion tree jointly.

\section{Randomized FIND}

The randomized FIND algorithm chooses a pivot, partitions the input, and
recurses only into the side containing the requested order statistic.  With
uniform random pivots, its expected number of comparisons is $O(n)$.

\begin{theorem}[Expected linear selection]
For every rank $k$, randomized FIND returns the $k$th smallest of $n$
distinct keys in expected $O(n)$ comparisons.
\end{theorem}

\begin{proof}[Proof sketch]
At each recursive call, the pivot rank is uniform in the current interval.
The expected size of the next interval is at most a constant fraction of the
current size after conditioning on whether the pivot falls in the middle
half.  Equivalently, charging the work at each level to surviving elements
gives a geometric series $O(n)+O(3n/4)+O((3/4)^2n)+\cdots=O(n)$.
\end{proof}

The implementation and the more detailed recurrence remain in
Chapter~\ref{ch:fundamentals}.

\section{Yao's minimax principle}

Let $D$ be a distribution over inputs and let $A$ be a randomized algorithm.
Yao's principle says that the expected cost of $A$ on the worst input is at
least the expected cost of the best deterministic algorithm under a suitable
hard input distribution:
\[
 \min_{A\text{ randomized}}\max_x\mathbb{E}[C(A,x)]
 \;\geq\;
 \max_D\min_{B\text{ deterministic}}\mathbb{E}_{x\sim D}[C(B,x)].
\]

\begin{proof}[Proof]
Fix a randomized algorithm $A$, viewed as a distribution over deterministic
algorithms $B$.  For every distribution $D$,
\[
 \max_x\mathbb{E}_A C(A,x)
 \geq \mathbb{E}_{x\sim D}\mathbb{E}_A C(A,x)
 =\mathbb{E}_A\mathbb{E}_{x\sim D}C(A,x)
 \geq\min_B\mathbb{E}_{x\sim D}C(B,x).
\]
Taking the maximum over $D$ proves the claim.
\end{proof}

To use the principle, students must specify the input distribution, prove a
lower bound for every deterministic algorithm, and then state exactly which
randomized guarantee follows.  The principle does not automatically prove a
worst-case lower bound for every cost notion without these quantifiers.

\section{Amplification}

Suppose a Monte Carlo decision algorithm is correct with probability at least
$1/2+\gamma$ for every fixed input, where $\gamma>0$.  Run independent copies
and return the majority answer.  Hoeffding's inequality gives error at most
$\exp(-2r\gamma^2)$ after $r$ copies.  Thus
$r\geq(\ln(1/\delta))/(2\gamma^2)$ gives error at most $\delta$.

For one-sided error, repeat until a desired confidence level is reached or
use the median of independent estimates.  Repetition must use fresh random
bits; reusing the same random seed does not amplify anything.

\section{Explicit lower-bound examples}

\subsection{Sorting}

For comparison sorting, choose the uniform distribution over all $n!$
relative orders.  A deterministic comparison tree of height $h$ has at
most $2^h$ leaves, so its expected depth under the uniform distribution is
$\Omega(n\log n)$.  Yao's principle transfers this to every Las Vegas
comparison sort on some input.

\subsection{Randomized search}

For searching an unsorted array for a marked position, choose the marked
position uniformly.  Every deterministic algorithm has expected
$(n+1)/2$ probes under this distribution.  Therefore no randomized algorithm
can have worst-case expected fewer than $(n+1)/2$ probes.

\subsection{Game-tree evaluation}

The game-tree lower-bound arguments in Chapter~\ref{ch:game-theoretic}
provide a stronger example: an input distribution can force every
deterministic decision strategy to inspect many leaves, even though a
randomized strategy performs substantially better on every fixed input.

\section{Student checklist}

For each randomized algorithm, students should answer:

\begin{enumerate}
\item Is it Las Vegas or Monte Carlo?
\item What is random: the input, the algorithm, or both?
\item Is the guarantee pointwise for every input or average-case?
\item What is the expected and high-probability runtime?
\item How is error amplified, and are the repetitions independent?
\item If a lower bound is claimed, what distribution and quantifiers does
      Yao's principle use?
\end{enumerate}

\begin{problem}[Quicksort indicators]
Derive the probability that two fixed ranks are compared by randomized
Quicksort and sum the resulting indicators.
\end{problem}

\begin{problem}[Yao distribution]
Give an input distribution proving an $\Omega(n)$ expected-probe lower bound
for randomized search in an unsorted array.
\end{problem}

\begin{problem}[Amplification]
An algorithm succeeds with probability $0.6$.  How many independent runs are
needed to reduce the majority-vote error below $10^{-9}$ using a Hoeffding
bound?
\end{problem}

\begin{problem}[Implementation]
Implement randomized Quicksort and randomized FIND.  Compare empirical
comparison counts with their expected bounds on sorted, reverse-sorted, and
random inputs.
\end{problem}
